\chapter*{Abstract}
\addcontentsline{toc}{chapter}{Abstract}
Missing or degraded electroencephalography (EEG) channels are a recurring issue in clinical studies, cognitive neuroscience experiments, and brain-computer interfaces. Sensor loss distorts the spatial topography of brain activity, affects event-related potentials, and may bias subsequent spectral analyses. This thesis studies EEG channel reconstruction under the framework of Graph Signal Processing (GSP), comparing classical geometric interpolation, instantaneous graph-based methods, and temporally regularized graph methods.

The central contribution is methodological and experimental: a reproducible, multi-dataset, multi-metric benchmark is built in which all evaluated methods, including classical baselines, undergo systematic hyperparameter optimization. The evaluation includes random and spatially contiguous missing-channel scenarios, mild-to-severe degradation levels, and amplitude, temporal-morphology, and spectral-fidelity metrics. The results indicate that temporal regularization, mainly represented by TRSS, provides consistent advantages in morphological and spectral fidelity under the evaluated protocol, especially when spatial information is sparse or locally degraded. At the same time, the thesis emphasizes that fair comparison narrows the apparent gap between modern algorithms and classical baselines; therefore, performance claims must be expressed with quantitative traceability and methodological caution.

\textbf{Keywords:} EEG, channel interpolation, graph signal processing, GSP, TRSS, Optuna, reconstruction, temporal regularization, spectral metrics.
